% =============================================================
% chapter_displaced.tex — кейс переміщених закладів П(ПТ)О
% Луганщини як спеціальний контекст моніторингу.
% Status: stub — наповнюється у Phase 1 / T1.5.
% =============================================================

\section{\MakeUppercase{Кейс переміщених закладів професійної (професійно-технічної) освіти Луганщини}}
\label{sec:displaced}

\subsection{Хронологія евакуації та фактичне розміщення}
\label{sec:displaced-history}

% TODO: спертися на дані МОН / ОДА / ETF про релокацію ЗП(ПТ)О Луганщини
% з 2014 та 2022 рр.; розрив юридичної vs. фактичної адреси.

\subsection{Виклики моніторингу контингенту й матеріально-технічної бази}
\label{sec:displaced-monitoring}

% TODO: облік контингенту, що фізично роззосереджений; ВПО серед здобувачів;
% евакуйована МТБ; дистанційний навчальний процес у тилу.

\subsection{Як ці виклики впливають на дизайн системи моніторингу}
\label{sec:displaced-design-impact}

% TODO: пов'язати з Частиною 2 — поле «фактична адреса», статус
% «евакуйований», ABAC-полісі для переміщених, спеціальні форми звітності.
